% block: matrix-fractional-power-heading page:220 type:section
\subsection{{Matrix Fractional Power}}\label{sec:matrix-fractional-power}

% block: matrix-fractional-power-definition page:220 type:ordinary-paragraph
{Matrix fractional power can be defined by using $\log(A)$ as follows:}
\begin{equation}\label{eq:matrix-fractional-power-definition}A^\alpha=\exp(\alpha\log(A)),\end{equation}
{where $0<\alpha<1$. The condition $0<\alpha<1$ looks to be too restrictive, but it is satisfactory in practice. In fact, if one wants to compute $A^{2.3}$, then $A^{2.3}$ can be decomposed by $A^{2.3}=A^2A^{0.3}$. Therefore, the problem is how to compute $A^{0.3}$ that corresponds to \eqref{eq:matrix-fractional-power-definition} with $\alpha=0.3$.}

% block: matrix-fractional-power-integral page:220 type:ordinary-paragraph
{Computing $A^\alpha$ via the definition in \eqref{eq:matrix-fractional-power-definition} requires two matrix functions: a matrix logarithm and a matrix exponential function. On the other hand, these matrix functions do not appear in the following integral form:}
\begin{equation}\label{eq:matrix-fractional-power-integral}A^\alpha=\frac{\sin(\alpha\pi)}{\alpha\pi}A\int_0^\infty(t^{1/\alpha}I+A)^{-1}\,\mathrm{d}t\quad(0<\alpha<1).\end{equation}

% block: matrix-fractional-power-de-formula page:220 type:ordinary-paragraph
{Application of the DE formula to \eqref{eq:matrix-fractional-power-integral} is considered in \BookCite{186} as described next:}
\begin{equation}\label{eq:matrix-fractional-power-de-integral}A^\alpha=\int_{-\infty}^\infty F_{\mathrm{DE}}(x)\,\mathrm{d}x,\end{equation}
{where}
\begin{equation}\label{eq:matrix-fractional-power-de-integrand}F_{\mathrm{DE}}(x)=t'(x)F(t(x)),\quad F(t)=\frac{\sin(\alpha\pi)}{\alpha\pi}A(t^{1/\alpha}I+A)^{-1}.\end{equation}

% block: matrix-fractional-power-algorithm-transition page:220 type:ordinary-paragraph
{The algorithm for computing $A^\alpha$ by the DE formula is described in Algorithm~\ref{alg:matrix-fractional-power-de}. For the theoretical details of Algorithm~\ref{alg:matrix-fractional-power-de}, see \BookCite{186}.}

% block: matrix-fractional-power-de-algorithm page:220 type:ordinary-paragraph
\begin{algorithm}[H]\caption{{$m$-point DE formula for computing $A^\alpha$}}\label{alg:matrix-fractional-power-de}\begin{algorithmic}[1]
\State \textbf{Input} $A\in\mathbb{R}^{n\times n}$, $\alpha\in(0,1)$, $\epsilon>0$, $m$
\State $l,r=\mathtt{GetInterval}(A,\alpha,\epsilon)$
\State {Set} $\widetilde{F}_{\mathrm{DE}}(x):=\exp(\alpha\pi\sinh(x)/2)\cosh(x)[\exp(\pi\sinh(x)/2)I+A]^{-1}$.
\State $h=(r-l)/(m-1)$
\State $T=h[\widetilde{F}_{\mathrm{DE}}(l)+\widetilde{F}_{\mathrm{DE}}(r)]/2+h\sum_{k=1}^{m-2}\widetilde{F}_{\mathrm{DE}}(l+kh)$
\State \textbf{Output} $\sin(\alpha\pi)AT/2\approx A^\alpha$
\State{}
\State \textbf{function} $\mathtt{GetInterval}(A,\alpha,\epsilon)$
\State {Compute $\|A\|$, $\|A^{-1}\|$.}
\State $a_1=[\alpha\pi(1+\alpha)\epsilon]/[4\sin(\alpha\pi)(1+2\alpha)]$, $a_2=(2\|A^{-1}\|)^{-\alpha}$
\State $a=\min\{a_1,a_2\}$
\State $\begin{aligned}b_1&=[\pi(1-\alpha)(2-\alpha)\epsilon]^{\alpha/(\alpha-1)}/[4\sin(\alpha\pi)(3-2\alpha)\|A\|]^{\alpha/(\alpha-1)},\\b_2&=(2\|A\|)^\alpha\end{aligned}$
\State $b=\max\{b_1,b_2\}$
\State $l=\operatorname{asinh}(2\log(a)/\alpha\pi)$, $r=\operatorname{asinh}(2\log(b)/\alpha\pi)$
\State \textbf{return} $l,r$
\State \textbf{end function}
\end{algorithmic}\end{algorithm}

% block: matrix-fractional-power-column-systems page:221 type:ordinary-paragraph
{From Algorithm~\ref{alg:matrix-fractional-power-de}, the $j$th column of computed $A^\alpha$ is $(\sin(\alpha\pi)AT/2)\boldsymbol{e}_j$. Therefore, from line 5 the $\widetilde{F}_{\mathrm{DE}}(l+ih)\boldsymbol{e}_j$'s must be computed using the following equations:}
\[e^{\alpha\pi\sinh(l+ih)/2}\cosh(l+ih)(e^{\pi\sinh(l+ih)/2}I+A)\boldsymbol{x}^{(i)}=\boldsymbol{e}_j\quad\text{for }i=0,1,\ldots,m-1.\]
{Similar to \eqref{eq:matrix-log-quadrature-systems} and \eqref{eq:matrix-log-shifted-systems}, the equations can be rewritten as shifted linear systems. Thus, if matrix $A$ is large and sparse, using (shifted) Krylov subspace methods will be a method of choice.}
